\section{集合的乘积}

\begin{enumerate}
\def\labelenumi{\arabic{enumi}.}
\tightlist
\item
  设 E, F 是两个集合，它们可以相同也可以不同。有序对 \((x, y)\) ，其中第一个元素 x 是 E 的任意元素，第二个元素 y 是 F 的任意元素，它们构成一个新集合的元素，该集合称为 E 与 F 的乘积，记为 \(E \times F\) ；E 和 F 称为 \(E \times F\) 的因子。两个有序对被认为相同，仅当它们具有相同的第一个元素和相同的第二个元素；换言之，关系“ \((x, y) = (x', y')\) ''等价于关系`` \(x = x'\) 并且 \(y = y''\) ”。若 z 是 \(E \times F\) 的任意元素，则关系“x 是有序对 \(z''\) 的第一个元素”是 x 的一个函数关系；它确定了一个从 \(E \times F\) 到 E 的映射，称为第一坐标函数，或第一投影，记为 \(pr_{1}\) 。除了说“x 是有序对 \(z''\) 的第一个元素”之外，我们也说“x 是 \(z''\) 的第一坐标（或投影）”或“ \(x = \mathrm{pr}_1z''\) ”。类似地，我们定义第二坐标函数，或第二投影，它是从 \(E \times F\) 到 \(F\) 上的映射，记为 \(\mathrm{pr}_2\) .
\end{enumerate}

关系“ \(x = \text{pr}_1 z\) 并且 \(y = \text{pr}_2 z\) ”等价于“ \(z = (x, y)\) ''.

函数 \(pr_{1}\) 到子集所成集合的扩张，按照一般约定用相同的符号表示，并再次称为第一投影（这里我们不用“坐标”一词）；第二投影的扩张也类似。

\begin{enumerate}
\def\labelenumi{\arabic{enumi}.}
\setcounter{enumi}{1}
\tightlist
\item
  E 的生成元 x 与 F 的生成元 y 之间的关系 R 是有序对 \((x, y)\) 的一个性质，因此定义了乘积 \(E \times F\) 的一个子集，称为 R 的图。反之， \(E \times F\) 的每个子集 A 都是 x 与 y 之间关系“ \((x, y) \in A\) ”的图。
\end{enumerate}

设 A 是 E 的子集，B 是 F 的子集。我们用 \(A \times B\) 表示 \(E \times F\) 的由 x 与 y 之间关系“ \(x \in A\) 并且 \(y \in B\) ”定义的子集。

\begin{enumerate}
\def\labelenumi{\arabic{enumi}.}
\setcounter{enumi}{2}
\tightlist
\item
  在以下命题中，X 和 X' 表示 E 的任意子集，Y 和 Y' 表示 F 的任意子集，Z 表示 E × F 的任意子集。
\end{enumerate}

\begin{enumerate}
\def\labelenumi{(\alph{enumi})}
\item
  关系“X × Y = φ”等价于“X = φ 或 Y = φ”。
\item
  如果 \(X \times Y \neq \emptyset\) ，关系“ \(X \times Y \subset X' \times Y''\) ”等价于“ \(X \subset X'\) 并且 \(Y \subset Y''\) ''.
\item
  对于所有 \(\mathbf{X},\mathbf{X}^{\prime},\mathbf{Y}\) 我们拥有
\end{enumerate}

\[
(\mathbf {X} \times \mathbf {Y}) \cup (\mathbf {X} ^ {\prime} \times \mathbf {Y}) = (\mathbf {X} \cup \mathbf {X} ^ {\prime}) \times \mathbf {Y}.\tag{22}
\]

\begin{enumerate}
\def\labelenumi{(\alph{enumi})}
\setcounter{enumi}{3}
\tightlist
\item
  对于所有 \(\mathbf{X},\mathbf{X}^{\prime},\mathbf{Y},\mathbf{Y}^{\prime}\) 我们拥有
\end{enumerate}

\[
(\mathbf {X} \times \mathbf {Y}) \cap (\mathbf {X} ^ {\prime} \times \mathbf {Y} ^ {\prime}) = (\mathbf {X} \cap \mathbf {X} ^ {\prime}) \times (\mathbf {Y} \cap \mathbf {Y} ^ {\prime}).
\]

\begin{enumerate}
\def\labelenumi{(\alph{enumi})}
\setcounter{enumi}{4}
\tightlist
\item
  对所有 X, Y，我们有
\end{enumerate}

\[
\mathrm{pr} _ {1} ^ {- 1} (\mathbf {X}) = \mathbf {X} \times \mathbf {F}, \quad \mathrm{pr} _ {2} ^ {- 1} (\mathbf {Y}) = \mathbf {E} \times \mathbf {Y}.\tag{24}
\]

\begin{enumerate}
\def\labelenumi{(\alph{enumi})}
\setcounter{enumi}{5}
\tightlist
\item
  如果 \(\mathbf{Y} \neq \emptyset\) ，那么对所有 \(\mathbf{X}\) 我们拥有
\end{enumerate}

\[
\operatorname{pr} _ {1} (\mathbf {X} \times \mathbf {Y}) = \mathbf {X}.\tag{25}
\]

\begin{enumerate}
\def\labelenumi{(\alph{enumi})}
\setcounter{enumi}{6}
\tightlist
\item
  对所有 Z，我们有
\end{enumerate}

\[
Z \subset \operatorname{pr} _ {1} (Z) \times \operatorname{pr} _ {2} (Z).\tag{26}
\]

\begin{enumerate}
\def\labelenumi{(\alph{enumi})}
\setcounter{enumi}{7}
\tightlist
\item
  令 \(a\) 为 \(\mathbf{E}\) 的一个元素。则映射 \((a, y) \to y\) （集合 \(\{a\} \times \mathbf{F}\) 到 \(\mathbf{F}\) 上的，即 \(\mathbf{pr}_2\) 到子集 \(\{a\} \times \mathbf{F}\) 的限制）是一对一的。
\end{enumerate}

\begin{enumerate}
\def\labelenumi{\arabic{enumi}.}
\setcounter{enumi}{3}
\item
  该映射

\[
(x, y) \rightarrow (y, x)\tag{27}
\]

是 \(E \times F\) 到 \(F \times E\) 上的一一映射，并称为典范映射。当 E 和 F 是同一个集合时，映射 (27) 称为典范对称；它此时是对合的。元素 \((x, y)\) （属于 \(E \times E\) ）在此对称下保持不变的，正是那些具有性质 x = y 的元素；这些元素组成的集合 \(\Delta\) 称为 \(E \times E\) 的对角线. 映射 \(x \to (x, x)\) 是 E 到 \(\Delta\) 上的双射，称为 E 到 \(E \times E\) 中的对角映射.

如果 Z 是 \(E \times F\) 的任意子集，Z 在 \(E \times F\) 到 \(F \times E\) 上的典范映射下的像记作 \(\overline{Z}^{1}\) . 如果 X 是 E 的任意子集且 Y 是 F 的任意子集，则

\[
\overline{\mathbf {X} \times \mathbf {Y}}^{1} = \mathbf {Y} \times \mathbf {X}.
\]

如果 x 与 y 之间的关系 R，作为有序对 \((x, y)\) 的一个性质，定义了 \(E \times F\) 的一个子集 A，那么同一关系，作为该对的一个性质 \((y, x)\) ，定义子集 \(\overline{A}^{1}\) 的 \(F \times E\) 。R 等价于以下每个关系 \((x, y) \in A\) , \((y, x) \in \overline{A}^{1}\) 。如果 E 和 F 是同一个集合，则当 \(A = \overline{A}^{1}\) 时，关系 R 及相应的子集 A 被称为对称的。对角线 \(\Delta\) （由相等关系定义的）是对称的。如果 Z 是 \(E \times E\) ，那么 \(Z \cup \overline{Z}^{1}\) 并且 \(Z \cap \overline{Z}^{1}\) 是对称的。

\end{enumerate}

\begin{enumerate}
\def\labelenumi{\arabic{enumi}.}
\setcounter{enumi}{4}
\tightlist
\item
  设A是集合E的一个子集，且f是从A到集合F的一个映射。关系“ \(x \in A\) 并且 \(y = f(x)\) '' 在 E 的任意元素 x 与 F 的任意元素 y 之间定义了一个子集，该子集属于 \(E \times F\) 称为函数 f 的图像。～若 B 是 E 中包含 A 的子集，且 g 是 f 到 B 的延拓（§ 2，第 13 号），则 f 的图像包含于 g 的图像中。
\end{enumerate}

反之，设 C 是 \(\mathbf{E} \times \mathbf{F}\) 的一个子集，使得对于每个 \(x \in \mathbf{E}\) ，至多存在一个 \(y \in \mathbf{F}\) 使得 \((x, y) \in \mathbf{C}\) 。则关系 \((x, y) \in \mathbf{C}\) （在一般元素 \(x\) （属于集合 \(\operatorname{pr}_1(\mathbf{C})\) ）与一个一般元素 \(y\) （属于 \(\mathbf{F}\) ）之间）是 \(y\) 中的函数关系，并确定 \(\operatorname{pr}_1(\mathbf{C})\) 到 \(\mathbf{F}\) 中的一个映射，其图像为 \(\mathbf{C}\) .

子集 C 的集合， \(E \times F\) 其具有性质“对所有 \(x \in E\) 至多存在一个 \(y \in F\) 使得 \((x, y) \in C\) ''（一个集合，它是 \(\mathfrak{P}(\mathbf{E} \times \mathbf{F})\) 因此可以与将E的子集映射到F的映射集合建立一一对应关系。

让 \(f\) 是从 \(\mathbf{E}\) 到 \(\mathbf{F}\) 的单射映射，并且让 \(g\) 是 \(f\) （视为 \(\mathbf{E}\) 到 \(f(\mathbf{E})\) 上的双射）的逆。如果 \(\mathbf{C}\) 是 \(f\) 的图像，则 \(g\) 的图像是 \(\overline{\mathbf{C}}^1\) .

\begin{enumerate}
\def\labelenumi{\arabic{enumi}.}
\setcounter{enumi}{5}
\tightlist
\item
  如果 \(f\) 是从 \(\mathbf{E}\) 到 \(\mathbf{F}\) 的映射，以及 \(\mathbf{C}\) 是其在 \(\mathbf{E} \times \mathbf{F}\) 中的图像，关系“ \(y = f(x)\) ”等价于“ \((x, y) \in \mathbf{C}\) ''。关系“ \(y \in f(\mathbf{X})\) '' 等价于“存在 \(x\) 使得 \(x \in \mathbf{X}\) 并且 \((x, y) \in \mathbf{C}\) ''.
\end{enumerate}

现在设K是E×F的任意子集，X是E的任意子集。用K(X)表示F中由所有满足关系“存在x使得 \(x \in X\) 并且 \((x, y) \in K\) ''”的元素y组成的子集；因此该关系等价于`` \(y \in K(X)\) ''。映射 \(X \to K(X)\) （从 \(\mathfrak{P}(E)\) 到 \(\mathfrak{P}(F)\) 的）称为由E × F的子集K所定义。注意，K(X)是集合 \(K \cap (X \times F)\) 的第二投影。如果 K 是映射 f 从 E 到 F 的图像，那么该映射 \(X \to K(X)\) 是 f 到子集集合的典范扩张。

\begin{enumerate}
\def\labelenumi{\arabic{enumi}.}
\setcounter{enumi}{6}
\tightlist
\item
  如果 \(x\) 是 \(\mathbf{E}\) 的典型元素，那么 \(x \to \mathrm{K}(\{x\})\) 是从 \(\mathbf{E}\) 到 \(\mathfrak{P}(\mathbf{F})\) 的映射，其值 \(\mathrm{K}(\{x\})\) （为简便起见也记为 \(\mathrm{K}(x)\) ）称为 \(\mathbf{K}\) 在 \(x\) 处的截面。关系 \((x, y) \in \mathbf{K}\) 等价于 \(y \in \mathrm{K}(x)\) .
\end{enumerate}

反之，每个映射 \(x \to \Phi(x)\) （从 \(E\) 到 \(\mathfrak{P}(F)\) ）可以以这种方式获得；因为该关系 \(y \in \Phi(x)\) 定义了子集 \(K\) （ \(E \times F\) 的），以及 \(\Phi(x)\) 正是 \(K\) 在 \(x\) 处的截面。集合 \(\mathfrak{P}(E \times F)\) 以及 \(E\) 到 \(\mathfrak{P}(F)\) 的映射的集合因此是一一对应的。

\begin{enumerate}
\def\labelenumi{\arabic{enumi}.}
\setcounter{enumi}{7}
\tightlist
\item
  每个映射 \(\mathbf{X} \to \mathbf{K}(\mathbf{X})\) （由子集 \(\mathbf{K}\) （ \(\mathbf{E} \times \mathbf{F}\) 的）所定义的）具有以下性质，这些性质推广了从 \(\mathbf{E}\) 到 \(\mathbf{F}\) 的映射的典范扩张的性质 (§2，第4和第5条)：
\end{enumerate}

\begin{enumerate}
\def\labelenumi{(\alph{enumi})}
\item
  \(\mathbf{K}(\phi)=\phi.\)
\item
  “X ⊂ Y” 蕴含 “K(X) ⊂ K(Y)”。
\item
  对所有 X, Y，我们有
\end{enumerate}

\begin{enumerate}
\def\labelenumi{(\arabic{enumi})}
\setcounter{enumi}{27}
\tightlist
\item
\end{enumerate}

\[
\mathbf {K} (\mathbf {X} \cup \mathbf {Y}) = \mathbf {K} (\mathbf {X}) \cup \mathbf {K} (\mathbf {Y}),\tag{29}
\]

\[
\mathbf {K} (\mathbf {X} \cap \mathbf {Y}) \subset \mathbf {K} (\mathbf {X}) \cap \mathbf {K} (\mathbf {Y}).
\]

若 K 且 \(K'\) 是 \(E \times F\) 使得 \(K \subset K'\) ，那么我们有 \(\mathrm{K}(\mathrm{X}) \subset \mathrm{K}'(\mathrm{X})\) 对于所有 \(X \subset E\) ；特别地， \(\mathrm{K}(x) \subset \mathrm{K}'(x)\) 对于所有 \(x \in E\) 。反之，若 \(\mathrm{K}(x) \subset \mathrm{K}'(x)\) 对于所有 \(x \in E\) ，那么 \(K \subset K'\) .

\begin{enumerate}
\def\labelenumi{\arabic{enumi}.}
\setcounter{enumi}{8}
\tightlist
\item
  E 的一般元素与 F 的一般元素之间的关系定义了 \(E \times F\) 的一个子集 K 以及 \(\overline{K}\) （ \(F \times E\) 的一个子集），因此有映射 \(X \to K(X)\) （从 \(\mathfrak{P}(E)\) 到 \(\mathfrak{P}(F)\) ）以及一个映射 \(Y \to \overline{K}(Y)\) （从 \(\mathfrak{P}(F)\) 到 \(\mathfrak{P}(E)\) ).
\end{enumerate}

如果 \(\mathbf{K}\) 是某个映射 \(f\) （从 \(\mathbf{E}\) 到 \(\mathbf{F}\) 的）的图像，映射 \(\mathbf{Y} \to \overline{\mathbf{K}}(\mathbf{Y})\) 是 \(f\) 的逆扩张.

应当注意，关系式(18)和(19)并不推广到映射 \(\mathbf{X} \to \mathbf{K}(\mathbf{X})\) 并且 \(\mathbf{Y} \to \overline{\mathbf{K}}^{1}(\mathbf{Y})\) ，当 \(\mathbf{K}\) 是 \(\mathbf{E} \times \mathbf{F}\) 的任意子集时.

\begin{enumerate}
\def\labelenumi{\arabic{enumi}.}
\setcounter{enumi}{9}
\tightlist
\item
  设E、F、G为三个集合，它们可能相同也可能不同，设A为 \(E \times F\) 的子集，并令 B 为 \(F \times G\) 的子集。则元素 \((x, z)\) （属于 \(E \times G\) ）中具有性质“存在 \(y \in F\) 使得 \((x, y) \in A\) 并且 \((y, z) \in B\) '' 的那些构成 \(E \times G\) 的一个子集，称为B与A的复合，记作 \(B \circ A\) ，或在无混淆风险时简称为BA。这里再次强调，复合的顺序至关重要。¶ 映射 \(X \to BA(X)\) （从 \(\mathfrak{P}(E)\) 到 \(\mathfrak{P}(G)\) ）是 \(Y \to B(Y)\) 与 \(X \to A(X)\) 的复合；换言之，对所有 \(X \subset E\) 我们拥有
\end{enumerate}

\[
\mathbf {B A} (\mathbf {X}) = \mathbf {B} (\mathbf {A} (\mathbf {X})).\tag{30}
\]

设H为另一个集合，不一定与E、F、G不同，并设C为 \(G \times H\) 的一个子集。于是我们有 \(\mathbf{C} \circ (\mathbf{B} \circ \mathbf{A}) = (\mathbf{C} \circ \mathbf{B}) \circ \mathbf{A}\) ；该集合也记作 \(C \circ B \circ A\) （或简记为CBA），并称为C、B、A按此顺序的复合。

让 \(f\) 是从 \(\mathbf{E}\) 到 \(\mathbf{F}\) 的映射，并且让 \(g\) 是从 \(\mathbf{F}\) 到 \(\mathbf{G}\) 的映射。如果 \(\mathbf{A}\) （相应地 \(\mathbf{B}\) ）是 \(f\) （相应地 \(g\) ）的图像，则复合 BA 是复合映射 \(g \circ f\) 的图像.

\begin{enumerate}
\def\labelenumi{\arabic{enumi}.}
\setcounter{enumi}{10}
\tightlist
\item
  我们有
\end{enumerate}

\[
\overline{\mathbf {B} \circ \mathbf {A}}^{1} = \overline {\mathbf {A}}^{1} \circ \overline {\mathbf {B}}^{1}.\tag{31}
\]

设 A, A' 是 E × F 的两个子集，设 B, B' 是 F × G 的两个子集。则关系

\[
" \mathbf {A} \subset \mathbf {A} ^ {\prime} \text {   且   } \mathbf {B} \subset \mathbf {B} ^ {\prime}" \quad \text { 蕴含 } \quad " \mathbf {B} \circ \mathbf {A} \subset \mathbf {B} ^ {\prime} \circ \mathbf {A} ^ {\prime}".
\]

设 A 是 \(E \times F\) 的子集， \(\Delta\) 为 \(E \times E\) 的对角线，以及 \(\Delta'\) 为 \(F \times F\) 的对角线。于是我们有

\[
\mathbf {A} \circ \Delta = \Delta^ {\prime} \circ \mathbf {A} = \mathbf {A}.\tag{32}
\]

\begin{enumerate}
\def\labelenumi{\arabic{enumi}.}
\setcounter{enumi}{11}
\tightlist
\item
  设E、F、G为三个集合，它们可能相同也可能不同。它们的积 \(E \times F \times G\) 是三元有序组的集合 \((x, y, z)\) ，其中 \(x \in E\) , \(y \in F\) ，以及 \(z \in G\) ，关系“ \((x, y, z) = (x', y', z')\) '' 等价于 `` \(x = x'\) 并且 \(y = y'\) 并且 \(z = z''\) ''。这三个映射 \((x, y, z) \to x\) , \((x, y, z) \to y\) , \((x, y, z) \to z\) （从 \(E \times F \times G\) 分别到E、F、G上）被称为第一、第二和第三坐标函数（或投影）；类似地，例如，指标为1、2的投影是映射
\end{enumerate}

\[
(x, y, z) \rightarrow (x, y)
\]

（ \(\mathbf{E} \times \mathbf{F} \times \mathbf{G}\) 到 \(\mathbf{E} \times \mathbf{F}\) 上的），并记作 \(\mathbf{pr}_{1,2}\) .

第2、3、4号的定义和命题很容易推广到三个集合的乘积。

此外，与其考虑三个集合的乘积 \(E \times F \times G\) ，我们同样可以考虑乘积 \((\mathbf{E} \times \mathbf{F}) \times \mathbf{G}\) ，它由对两个集合作乘积这一运算的两次应用而得到。事实上， \((x, y, z) \to ((x, y), z)\) 是 \(E \times F \times G\) 到 \((\mathbf{E} \times \mathbf{F}) \times \mathbf{G}\) 上的一个一一映射，称为典范的。类似地，我们定义 \(E \times F \times G\) 到集合 \(E \times (F \times G)\) 上的一一映射，以及到所有由 \(E \times F \times G\) , \((\mathbf{E} \times \mathbf{F}) \times \mathbf{G}\) ，以及 \(E \times (F \times G)\) 通过置换三个字母E、F、G而得到的集合上的一一映射。

对于三个以上集合的乘积，也有类似的定义和性质。

\begin{enumerate}
\def\labelenumi{\arabic{enumi}.}
\setcounter{enumi}{12}
\tightlist
\item
  如果一个函数 \(f\) ，其取值在任意集合 \(\mathbf{E}'\) 中，定义在三个集合的乘积 \(\mathbf{E}\) , \(\mathbf{F}\) , \(\mathbf{G}\) 上，则称它为三变量函数，其中每个变量分别遍历其中一个集合 \(\mathbf{E}\) , \(\mathbf{F}\) , \(\mathbf{G}\) 的值 \(f\) 在元素 \((x, y, z)\) 的 \(\mathbf{E} \times \mathbf{F} \times \mathbf{G}\) 处记作 \(f(x, y, z)\) .
\end{enumerate}

让 \(a\) 为 \(\mathbf{E}\) 的任意元素。则 \((y, z) \to f(a, y, z)\) 是从 \(\mathbf{F} \times \mathbf{G}\) 到 \(\mathbf{E}'\) 的映射，称为由 \(f\) 确定的部分映射（或函数），对应于该值 \(a\) （ \(x\) 所取的）；它也是 \(f\) 与映射 \((y, z) \to (a, y, z)\) （从 \(\mathbf{F} \times \mathbf{G}\) 到 \(\mathbf{E} \times \mathbf{F} \times \mathbf{G}\) ）的复合.

类似地，如果 \(b\) 是 \(\mathbf{F}\) 的一个元素，那么 \(z \to f(a, b, z)\) 是从 \(\mathbf{G}\) 到 \(\mathbf{E}'\) 的一个映射，称为由 \(f\) 确定的部分映射，对应于值 \(a, b\) （ \(x, y\) 所取的）.

反之，设g为从E到 \(E'\) 的映射。则 \((x, y, z) \to g(x)\) 是 \(E \times F \times G\) 到 \(E'\) 的一个映射 h，使得由 h 决定的、E 到 \(E'\) 的部分映射（对应于 y 和 z 的任何值）都与 g 相同。这一事实通常表述为：关于自变量 x 的函数总可以被视为关于在给定时刻需要考虑的所有自变量（当然包括 x）的函数。

\begin{enumerate}
\def\labelenumi{\arabic{enumi}.}
\setcounter{enumi}{13}
\tightlist
\item
  让 \(f, g, h\) 分别是从 \(\mathbf{E}\) 到 \(\mathbf{E}'\) , 从 \(\mathbf{F}\) 到 \(\mathbf{F}'\) , 从 \(\mathbf{G}\) 到 \(\mathbf{G}'\) 的三个映射。映射 \((x, y, z) \to (f(x), g(y), h(z))\) （从 \(\mathbf{E} \times \mathbf{F} \times \mathbf{G}\) 到 \(\mathbf{E}' \times \mathbf{F}' \times \mathbf{G}'\) ）记作 \((f, g, h)\) 或 \(f \times g \times h\) ，称为 \(f, g, h\) 到乘积的扩张。如果所有三个 \(f, g, h\) 是单射（相应地，满射、双射），则 \(f \times g \times h\) 是单射（相应地，满射、双射）。
\end{enumerate}

在本小节及上一小节中，我们仅考虑三个集合的情形，只是为了明确概念；对于任意有限个集合，类似的讨论同样成立。

